\documentclass[11pt]{article}
\usepackage{amsmath,amssymb,amsthm,hyperref}
\newtheorem{theorem}{Theorem}
\newtheorem{proposition}[theorem]{Proposition}
\newtheorem{lemma}[theorem]{Lemma}
\title{Shared Partial Fractions and Selector Flows}
\author{}
\date{24 November 2026}
\begin{document}
\maketitle

\paragraph{Status.}
These are complete written identities and interaction restrictions, not
Lean-checked results and not an asymptotic upper bound. The denominator
conditions of the original attention model are retained. In particular,
the explicit examples below have strictly positive slopes in every
original input coordinate and are positive on the whole cube.

\section{Which affine differences have affine reciprocals?}

Write the cube in centered coordinates $z_i\in\{-1,1\}$.

\begin{proposition}\label{prop:inverse}
Suppose the affine function $L=c+\sum_i a_i z_i$ never vanishes on the
cube. Its reciprocal is affine on the cube if and only if one of the
following holds:
\begin{enumerate}
\item $L$ depends on at most one coordinate;
\item $L=a_i z_i+a_j z_j$ for distinct $i,j$, where
      $a_i a_j\ne0$ and $a_i^2\ne a_j^2$.
\end{enumerate}
In the second case,
\[
 L^{-1}=\frac{a_i z_i-a_j z_j}{a_i^2-a_j^2}.
\]
\end{proposition}
\begin{proof}
Suppose $L(d+\sum_i b_i z_i)=1$. Comparing the distinct multilinear
coefficients gives
\[
 cb_i+da_i=0,\qquad a_i b_j+a_j b_i=0\quad(i\ne j).
\]
If $a_i=0$ and some $a_j\ne0$, the second equation gives $b_i=0$.
On the support of $a$, put $t_i=b_i/a_i$. Distinct supported indices
satisfy $t_i+t_j=0$. Three supported indices therefore force all
$t_i=0$, then $d=0$, contradicting the constant coefficient equation.
Thus the support has size at most two. If it has exactly two, write
$t_i=t$ and $t_j=-t$. The linear equations imply $ct+d=0$ and
$-ct+d=0$. If $c\ne0$, this again gives $t=d=0$, a contradiction.
Hence $c=d=0$, and the constant equation is
$t(a_i^2-a_j^2)=1$, giving the stated formula. A nonvanishing function
of one Boolean coordinate has an affine reciprocal by interpolation.
\end{proof}

\section{A general flow identity}

Let $G=(V,E)$ be a finite oriented graph. Give every vertex an admissible
affine denominator $D_v$. For an edge $e=(u,v)$ suppose the difference
$D_v-D_u$ is nonzero on the cube and has an affine reciprocal
\[
 R_e=r_{e0}+\sum_i r_{ei}x_i.
\]
Give each edge an arbitrary affine weight
$L_e=l_{e0}+\sum_i l_{ei}x_i$. Let $B_{ve}$ be $+1$ at the tail of
an edge, $-1$ at its head, and zero otherwise.

\begin{theorem}[Exact quadratic flow criterion]\label{thm:flow}
On the Boolean cube,
\[
 \sum_{e=(u,v)}\frac{L_e}{D_uD_v}
 =\sum_{v\in V}\frac{N_v}{D_v},\qquad
 N_v=\sum_e B_{ve}R_eL_e.
 \tag{1}
\]
The displayed numerators $N_v$, after Boolean reduction, are all affine
if and only if for every vertex $v$ and every pair $i<j$,
\[
 \sum_e B_{ve}(r_{ei}l_{ej}+r_{ej}l_{ei})=0.
 \tag{2}
\]
When (2) holds, the score in (1) uses at most $|V|$ admissible heads,
regardless of the number of edge products.
\end{theorem}
\begin{proof}
The identity $R_e(D_v-D_u)=1$ gives
$L_e/(D_uD_v)=R_eL_e/D_u-R_eL_e/D_v$. Sum over the edges.
The reduced product of two affine polynomials has degree at most two;
its coefficient of $x_i x_j$ is exactly
$r_{ei}l_{ej}+r_{ej}l_{ei}$. This proves necessity and sufficiency
for the displayed numerators to be affine. Their constant coefficients
are $\sum_e B_{ve}r_{e0}l_{e0}$, and their coefficients of $x_i$ are
\[
 \sum_e B_{ve}
 (r_{e0}l_{ei}+r_{ei}l_{e0}+r_{ei}l_{ei}).
\]
No denominator was changed. The criterion concerns this particular
partial-fraction construction; it does not assert uniqueness of all
possible representations using the same denominators.
\end{proof}

The theorem is algebraic for every allowed choice of the denominator
geometry, including geometry chosen after the target truth table.
It does not assume the denominators are fixed independently of that
table. Condition (2) says that every quadratic coefficient separately
forms a conserved edge flow.

\section{An explicit nonzero triangle}

Take any number of Boolean coordinates, one of which is $x$, and put
$D=1+\sum_i x_i$. Define
\[
 D_0=D+1+x,\qquad D_1=D+2+3x,\qquad D_2=D+4+6x.
\]
Every denominator is strictly positive and every slope is positive.
For every affine function $L$ of all the inputs,
\[
 -\frac{3L}{2D_0D_1}-\frac{10L}{3D_1D_2}
 +\frac{24L}{5D_2D_0}
 =\frac{L}{10D_0}-\frac{L}{6D_1}+\frac{L}{15D_2}.
 \tag{3}
\]
To verify this without division, multiply by $D_0D_1D_2$. The left
numerator minus the right numerator is $Lx(1-x)$, which vanishes
on the cube. The cancellation can also be read as a circulation.
For the oriented edges $01,12,20$, the affine reciprocals of the
denominator differences have constant coefficients
$1,1/2,-1/3$ and $x$ coefficients $-2/3,-3/10,5/24$.
The edge weights $-3/2,-10/3,24/5$ make their $x$ coefficients all
equal to one. Their divergence vanishes at each vertex, leaving the
constant residues $1/10,-1/6,1/15$ in (3).

\section{Several one-coordinate selectors}

Suppose each reciprocal $R_e$ depends on a single coordinate
$x_{c(e)}$, with a nonzero coefficient in that coordinate. Regard
$c(e)$ as the edge color. For $i<j$, let $G_{ij}$ consist of the edges
whose colors belong to $\{i,j\}$.

\begin{proposition}[Only bichromatic cycles carry quadratic cancellation]
For every $i<j$, equation (2) is a conserved scalar flow supported
on $G_{ij}$. If $G_{ij}$ is a forest, every edge contribution to the
coefficient of $x_i x_j$ is zero individually. More generally, that
coefficient flow is a linear combination of cycles in $G_{ij}$.
\end{proposition}
\begin{proof}
An edge of any other color has $r_{ei}=r_{ej}=0$. Thus the relevant
coefficient flow is supported on $G_{ij}$ and has zero divergence.
A nonzero flow on a finite forest would have a nonzero edge incident
to a leaf of its nonzero support, contradicting zero divergence there.
For the general statement, subtract a suitable multiple of a cycle
through a nonzero edge until fewer nonzero edges remain, and induct.
Equivalently, this is the usual cycle-space description of the kernel
of a graph's oriented incidence matrix.
\end{proof}

In particular, on a simple cycle using at least three different
coordinate colors, no quadratic terms can cancel across the whole
cycle. Each two-color subgraph is a proper subgraph of that cycle and
is a forest. This is an interaction restriction on this exact flow
mechanism, not on arbitrary attention scores.

\section{A genuine three-coordinate cycle}

Centered two-coordinate differences from Proposition~\ref{prop:inverse}
escape the preceding one-coordinate restriction. The following example
has a nonzero quadratic flow involving three coordinates.

Choose positive pairwise distinct $a,b,c$, and set $z_i=2x_i-1$ for
$i=1,2,3$. Let
\[
 D_1=T,\qquad D_2=T+az_1+bz_2,\qquad
 D_3=T+az_1+cz_3.
\]
Here $T=C+K\sum_i x_i$, where $K>0$ and $C>a+b+c$, makes all three
denominators positive with strictly positive slopes in every coordinate.
Put
\[
 \begin{split}
 L_{12}&=(a^2-b^2)(bz_2+cz_3),\\
 L_{23}&=(b^2-c^2)(-az_1+cz_3),\\
 L_{31}&=(a^2-c^2)(az_1-bz_2).
 \end{split}
\]
Then the exact identity
\[
 \frac{L_{12}}{D_1D_2}+\frac{L_{23}}{D_2D_3}
 +\frac{L_{31}}{D_3D_1}
 =\frac{a^2-b^2}{D_1}+\frac{b^2-c^2}{D_2}
 +\frac{c^2-a^2}{D_3}
 \tag{4}
\]
holds on the whole cube. Indeed the reciprocal differences are
\[
 R_{12}=\frac{az_1-bz_2}{a^2-b^2},\quad
 R_{23}=\frac{-bz_2-cz_3}{b^2-c^2},\quad
 R_{31}=\frac{-az_1+cz_3}{a^2-c^2}.
\]
If $Q=abz_1z_2+acz_1z_3-bcz_2z_3$, their products with the edge
weights are respectively $Q-b^2$, $Q-c^2$, and $Q-a^2$.
The common quadratic part cancels by circulation, yielding (4).
Multiplying every term of (4) by an arbitrary affine function $M$
still leaves exactly three affine-ratio heads on the right. Thus a
shared cancellation can remove higher-degree edge numerators jointly;
it need not remove them separately on each product term.

\section{Connected flow geometry cannot give a superconstant saving}

This restriction also permits a counting argument uniform over
target-dependent choices of the denominators. It applies only when the
affine-reciprocal-difference graph is connected.

\begin{lemma}[A two-block sign-pattern bound]\label{lem:count}
Let $F_1,\ldots,F_m$ be polynomials in blocks of $p$ and $q$ real
variables, of degrees at most one and $d$, respectively, where $d\ge1$.
Put $k=p+q$ and suppose $m\ge k\ge1$. There is an absolute constant
$A$ such that the number of their strict sign patterns is at most
\[
 \left(\frac{Am}{k}\right)^k d^q.
 \tag{5}
\]
\end{lemma}
\begin{proof}
Use Theorems 12 and 13 of Basu and Rizzie,
\emph{Multi-degree bounds on the Betti numbers of real varieties and
semi-algebraic sets and applications},
\href{https://www.math.purdue.edu/~sbasu/multidegree-may-24-2017.pdf}
{author's version of 24 May 2017}, page 11.
Use two blocks with degree bounds $2$ and $\max(2,d)$. Their quantity
$G_{\min}$ is bounded by $A_0^k d^q$ for an absolute $A_0$,
uniformly in the number of equations. Their sign-condition bound with
Betti index zero is therefore at most
\[
 A_0^k d^q\sum_{j=0}^k 4^j\binom mj
 \le (4eA_0m/k)^k d^q.
\]
The last inequality follows by setting $t=k/m\le1$ and bounding
$\sum_{j=0}^k\binom mj\le t^{-k}(1+t)^m\le(em/k)^k$.
Every realized strict sign condition contributes at least one connected
component, so its count is bounded by the cited Betti-number sum.
\end{proof}

\begin{theorem}\label{thm:connected}
Consider all scores on the $n$-cube consisting of a bias and $H$ affine
ratios with positive admissible denominators. Restrict this family by
requiring the following graph on the denominators to be connected:
two denominators are adjacent when their difference is nonvanishing
on the cube and has an affine reciprocal there. The denominators and
the graph may depend arbitrarily on the target function.
There is a constant $c>0$ such that, for all sufficiently large $n$,
realizing every Boolean sign pattern in this restricted family
requires $H\ge c2^n/n$.
\end{theorem}
\begin{proof}
Index the denominators and choose a spanning tree rooted at vertex one.
There are at most $H^{H-1}$ choices of rooted parent lists, including
an overcount of lists that are not trees. By
Proposition~\ref{prop:inverse}, every difference along a tree edge
has at most two nonzero affine coefficients in centered coordinates.
Overcount by assigning an ordered pair of indices from $0,\ldots,n$
to each edge, where index zero denotes the constant coefficient.
This gives at most $(n+1)^{2(H-1)}$ choices.

For each such discrete choice, specify the root denominator with
$n+1$ parameters and each edge difference with two. All denominator
coefficients are then linear functions of
$q=n+2H-1$ parameters. Allow all affine numerators and a bias,
using $p=H(n+1)+1$ further parameters. At a fixed cube vertex,
multiply the score by the positive product of all denominators.
The resulting polynomial has degree at most one in the numerator
block and at most $H$ in the denominator block. Put
$N=2^n$ and $k=p+q=H(n+3)+n$. If $k\le N$, Lemma~\ref{lem:count}
bounds the number of strict patterns in the entire restricted family by
\[
 H^{H-1}(n+1)^{2(H-1)}
 \left(\frac{AN}{k}\right)^k H^{n+2H-1}.
 \tag{6}
\]
Whole-cube positivity and slope restrictions only reduce this count.

To finish, suppose along some sequence $n\to\infty$ that
$Hn/N\to0$. Then $k/N\to0$ and
$k\log(AN/k)=o(N)$. Also $H\le N$ eventually, so
$(n+3H)\log H\le(n+3H)n\log2=o(N)$, and
$2H\log(n+1)=o(N)$. The logarithm of (6) is therefore $o(N)$,
strictly less than $N\log2$ eventually. Such a family cannot contain
all $2^N$ strict sign patterns. If the stated uniform constant $c$
did not exist, one could select an increasing sequence of dimensions
with universal budgets satisfying $Hn/N\to0$, giving exactly this
contradiction. The case $H=0$ supplies only constant patterns.
\end{proof}

Thus the target-dependent geometry in a single connected flow network
does not rescue this particular construction. The restriction is
substantive: general pairs of admissible affine denominators need not
have affine-reciprocal differences. Non-affine edge reciprocals, or a
different cancellation identity, remain outside Theorem~\ref{thm:connected}.

\section{Scope of the remaining problem}

The identities provide exact assembly rules with original admissible
denominators. They do not provide an improved universal head bound.
In particular,
counting all edge weights as independent output information is invalid:
the conserved quadratic parts disappear from the score. A successful
extension must exploit target-dependent choices of denominators or
graph geometry while exhibiting the actual surviving strict cube signs.
The bichromatic-cycle restriction is uniform over such choices for
one-coordinate reciprocal differences, but the three-coordinate
example shows that it cannot be extended unchanged to all affine
reciprocal differences.

\end{document}
